\section{Heights and coordinate tests}\label{sec:heights}

We apply the prediction device to the incidence pieces of
\Cref{inc:piece}. Their coordinates are real algebraic numbers.
The first task is to measure the variation of edge directions simultaneously
at all absolute values of a number field. We will prove that this variation
cannot remain bounded. We then construct coordinate labels whose agreement
records a threshold comparison, with a bounded total error.

Write
\[
 z_x=x_1+i x_2,\qquad w_x=x_1-i x_2
\]
for a point or center $x=(x_1,x_2)$. For every edge $e=(p,q)$, put
\begin{equation}\label{ht:jumps}
 d_e=z_p-z_q,\qquad w_p-w_q=d_e^{-1}.
\end{equation}
The second equality follows from the unit-distance equation, and in
particular $d_e\ne0$. Distinct real points have distinct $z$-coordinates and
distinct $w$-coordinates. These nonzero differences remain nonzero under
every field embedding. At an embedding other than the original one,
we use the two embedded coordinates separately; they need not be complex
conjugates.

\subsection{A weighted product formula}

For each configuration choose a finite normal number field $\mathcal N$
containing all its coordinates and $i$, and let
$d_{\mathcal N}=[\mathcal N:\Q]$.
The following construction allows repeated absolute values; the repetitions
make their weights particularly simple.

\begin{lemma}[Weighted absolute values]\label{ht:product}
There is a countable collection of absolute-value entries $v$ on
$\mathcal N$, each of weight $1/d_{\mathcal N}$, with the following properties.
Its archimedean entries are $|u|_v=|\sigma(u)|$ for the $d_{\mathcal N}$ embeddings
$\sigma:\mathcal N\to\C$, so their total weight is one. All other entries
are nonarchimedean. For each nonzero $u\in\mathcal N$, only finitely many
nonarchimedean entries have $|u|_v\ne1$, and
\begin{equation}\label{ht:product-formula}
 \sum_v^*\log|u|_v=0.
\end{equation}
Here and below a star denotes the sum with the stated weights.
\end{lemma}

\begin{proof}
Use the classification of number-field places and the product formula
\cite[Theorems~7.14--7.15]{ProductFormula}. At a finite place
\(\mathfrak p\) above a rational prime \(\ell\), normalize the genuine
absolute value by \(|\ell|_{\mathfrak p}=\ell^{-1}\), and put
\(n_{\mathfrak p}=[\mathcal N_{\mathfrak p}:\Q_\ell]\), where
\(\mathcal N_{\mathfrak p}\) is the completion. Use
\(n_{\mathfrak p}\) copies of this absolute value, each of weight
\(1/d_{\mathcal N}\). Milne's local normalization formula
\cite[Lemma~8.6]{ProductFormula} converts his finite-place value to
\(|\cdot|_{\mathfrak p}^{n_{\mathfrak p}}\). His complex-place value
is squared modulus, so we instead use ordinary modulus once for each
embedding, counting conjugate embeddings separately. Taking logarithms
of the product formula and dividing by \(d_{\mathcal N}\) proves
\eqref{ht:product-formula}. There are exactly \(d_{\mathcal N}\)
embeddings, giving archimedean weight one. Finally, the principal
fractional ideal of any nonzero \(u\) has finitely many prime factors
\cite[Theorem~3.20]{ProductFormula}; at every other finite place
\(|u|_v=1\).
\end{proof}

Define $\log^+r=\max(0,\log r)$ for $r>0$ and $\log^+0=0$.
Set $h(0)=0$, and put
\[
 h(u)=\sum_v^*\log^+|u|_v\qquad(u\ne0).
\]

\begin{lemma}[Height inequalities]\label{ht:height}
The value $h(u)$ is independent of the chosen finite normal field and of
the extensions of the prime absolute values in \Cref{ht:product}.
For $u\in\mathcal N^\times$,
\begin{equation}\label{ht:height-inversion}
 h(u^{-1})=h(u),\qquad
 \sum_v^*|\log|u|_v|=2h(u).
\end{equation}
For arbitrary $u,v\in\mathcal N$,
\begin{equation}\label{ht:height-inequalities}
 h(uv)\le h(u)+h(v),\qquad
 h(u+v)\le h(u)+h(v)+\log2.
\end{equation}
If a random $X\in\mathcal N$ is independent of a uniformly sampled
embedding $\sigma:\mathcal N\to\C$, then, for $B>1$ and $0<r<1$,
\begin{align}
 \Prob\{|\sigma(X)|>B\}&\le\frac{\E h(X)}{\log B},
       \label{ht:large-tail}\\
 \Prob\{X\ne0,\ |\sigma(X)|<r\}
      &\le\frac{\E h(X)}{\log(1/r)}.
       \label{ht:small-tail}
\end{align}
\end{lemma}

\begin{proof}
For a finite extension \(L/\mathcal N\), the decomposition into
completions \cite[Proposition~8.2]{ProductFormula} gives
\(\sum_{\mathfrak q\mid\mathfrak p}[L_{\mathfrak q}:\mathcal N_{\mathfrak p}]
=[L:\mathcal N]\). Hence
\[
 \sum_{\mathfrak q\mid\mathfrak p}
 \frac{[L_{\mathfrak q}:\Q_\ell]}{[L:\Q]}
 =\frac{[\mathcal N_{\mathfrak p}:\Q_\ell]}{[\mathcal N:\Q]}.
\]
The extended absolute values restrict to the chosen normalization on
\(\mathcal N\), so the finite-place contribution of any
\(u\in\mathcal N\) is unchanged. Each embedding of \(\mathcal N\)
has \([L:\mathcal N]\) extensions, giving the same conclusion at
infinity. A common finite normal overfield compares any two field
choices. The place list includes all normalized extensions with their
local-degree multiplicities, so their enumeration involves no further
choice. This proves height independence; \(h(0)=0\) is fixed.

The product formula equates the sums of positive and negative logarithms,
giving \eqref{ht:height-inversion}. Multiplication gives the first height
inequality termwise. The triangle inequality contributes at most $\log2$
at an archimedean entry and no constant at a nonarchimedean entry; the
total archimedean weight is one. This proves the second inequality,
including zero summands by the stated convention. For the tail estimates,
the average positive logarithm over embeddings is at most $h(X)$.
For $X\ne0$, its average negative logarithm is at most $h(X^{-1})=h(X)$.
Markov's inequality proves both assertions.
\end{proof}

\subsection{Thresholds and their total variation}\label{sec:levels}

For every entry $v$ define
\[
 n_e(v)=-\log|d_e|_v.
\]
Choose a finite closed interval $I_v=[\ell_v,u_v]$ containing $0$ and all
$n_e(v)$, enlarged so that for distinct $p,p'\in P$,
\begin{equation}\label{ht:range}
 \ell_v\le-\log|z_p-z_{p'}|_v,\qquad
 u_v\ge\log|w_p-w_{p'}|_v.
\end{equation}
All the differences here are nonzero. By \Cref{ht:product}, we can take
$I_v=[0,0]$ at all but finitely many entries. Define the level space
$\mathcal L$ to be the disjoint union of the $\{v\}\times I_v$, with
weighted Lebesgue measure:
\[
 \int_{\mathcal L}g
   :=\sum_v^*\int_{\ell_v}^{u_v}g(v,s)\dd s.
\]
We usually suppress the domain $\mathcal L$ in an integral.
At a level $(v,s)$ put
\begin{equation}\label{ht:levels}
 \begin{aligned}
 Y_e(v,s)&=\mathbf1_{\{n_e(v)>s\}},&
 \pi(v,s)&=\Prob_E\{Y_e(v,s)=1\},\\
 f(v,s)&=\min(\pi(v,s),1-\pi(v,s)),& S&=\int f.
 \end{aligned}
\end{equation}
Enlarging any $I_v$ further does not change $S$, because all edge bits
agree outside the range of the $n_e(v)$.

The choice of threshold turns distance between logarithmic profiles into
bit disagreement. For independent uniform edges $e,e'$, Tonelli's theorem
and \eqref{ht:height-inversion} give
\begin{equation}\label{ht:independent-profiles}
 2\E_{e,e'}h(d_e/d_{e'})
 =\E_{e,e'}\sum_v^*|n_e(v)-n_{e'}(v)|
 =2\int\pi(1-\pi)\le2S.
\end{equation}

\subsection{Excluding bounded total variation}

We first use matching labels to transfer height control from edges to
points. The argument is useful because it bounds heights of differences
before any translation has been chosen.

\begin{lemma}[Transferring edge heights]\label{ht:transfer}
Suppose a probability law $\omega$ on two finite vertex sets $P\times Q$,
with point marginal $\mu$, has the following property: for all assignments of finite
labels to its vertices, two independent $\mu$-points disagree with
probability at most $C_0$ times the edge disagreement probability.
For any coordinate $U:P\cup Q\to\mathcal N$,
\[
 \E_{p,p'\sim\mu}h(U_p-U_{p'})
   \le C_1\bigl(\E_{(p,q)\sim\omega}h(U_p-U_q)+1\bigr),
\]
where $C_1$ depends only on $C_0$.
\end{lemma}

\begin{proof}
Fix an entry $v$ and $s\ge0$. At a nonarchimedean entry give two vertices
the same label if their $U$-coordinates differ by at most $e^s$ in
absolute value. The strong triangle inequality makes this an equivalence
relation. At an archimedean entry, use the cell of a square grid in
$\C$ of side $e^s$, shifted by a uniformly random fractional vector.
Only the labels occurring on the finite vertex set matter.

For complex numbers whose distance is $a$, the probability of different
grid cells is at most
\[
 \min(1,\sqrt2\,a e^{-s}).
\]
Indeed a mismatch in either real coordinate has probability at most its
coordinate distance divided by the side length, and the two coordinate
distances sum to at most $\sqrt2 a$. Its integral over $s\ge0$ is at most
$C(1+\log^+ a)$. At a nonarchimedean entry the integral of mismatch is
exactly $\log^+|U_p-U_q|_v$, with value zero for a zero difference.
Thus the integrated edge mismatch, averaged over shifts at infinity,
is at most $C(\E_{(p,q)\sim\omega}h(U_p-U_q)+1)$.

Apply the assumed label inequality at each fixed level and shift.
For a point pair at a nonarchimedean entry, mismatch is exactly the
event $|U_p-U_{p'}|_v>e^s$. At infinity, distance greater than
$\sqrt2 e^s$ forces mismatch for every shift. Integration of these
tail events bounds the point-pair height, with an additional
$\log\sqrt2$ from the total archimedean weight. This proves the result.
\end{proof}

\begin{proposition}[The bounded alternative is impossible]\label{ht:bounded}
For a sequence of the incidence pieces in \Cref{inc:piece}, with the
level spaces and quantities $S$ in \eqref{ht:levels}, no subsequence
has $S$ bounded. Consequently $S\to\infty$.
\end{proposition}

\begin{proof}
Suppose $S\le S_0$ on a subsequence. All constants in this proof may
depend on the fixed $S_0$. By \eqref{ht:independent-profiles}, there is an
edge value $d_0$ such that
\[
 \E_Eh(d_e/d_0)\le S_0.
\]
Temporarily replace the two coordinates by $z/d_0$ and $d_0w$.
Their edge jumps are $D_e=d_e/d_0$ and $D_e^{-1}$, whose mean heights
are at most $S_0$. \Cref{inc:matching} supplies the label hypothesis of
\Cref{ht:transfer}. Applying it to both coordinates bounds the sum of
the two mean point-pair difference heights. Averaging over the second
point then chooses $p_0\in P$ such that the translated coordinates
\[
 Z_x=(z_x-z_{p_0})/d_0,\qquad W_x=d_0(w_x-w_{p_0})
\]
satisfy a fixed bound on $\E_\mu(h(Z_p)+h(W_p))$.
Since $Z_q=Z_p-D_{pq}$ and $W_q=W_p-D_{pq}^{-1}$ on an edge,
\Cref{ht:height} also bounds the center means. Thus there is a fixed
$C_h<\infty$ with
\begin{equation}\label{ht:normalized-means}
 \E_\mu(h(Z_p)+h(W_p))
 +\E_\kappa(h(Z_q)+h(W_q))
 +\E_E(h(D_e)+h(D_e^{-1}))\le C_h.
\end{equation}
The edge identity is still
$(Z_p-Z_q)(W_p-W_q)=1$.

We next give the selection argument, keeping the embedding independent
of the points it tests. For \emph{any} fixed finite history map $H$ on
$P$, use the joint law
\begin{equation}\label{ht:coupling}
 \mu(p)\,\mu(\widehat p\mid H(\widehat p)=H(p))
             \prod_{j=1}^3\Prob(q_j\mid p),
\end{equation}
and independently choose a uniformly random embedding $\sigma$.
Both $p$ and $\widehat p$ have marginal $\mu$; each $q_j$ has marginal
$\kappa$; and each $(p,q_j)$ has the original edge law.
In particular, the three jumps $D_j=D_{pq_j}$ have the correct individual
edge marginals. No assertion of mutual independence of these jumps is
needed. The height inequality bounds
$\E h(Z_{\widehat p}-Z_p)$ and each $\E h(D_j-D_l)$ using only
\eqref{ht:normalized-means}. These bounds are uniform in $H$.

Choose a fixed $B\ge2$ so large that, by \eqref{ht:large-tail}, the
probability that any of
\[
 Z_p,W_p,Z_{\widehat p},W_{\widehat p},
 (Z_{q_j},W_{q_j},D_j,D_j^{-1})_{j=1}^3
\]
has embedded absolute value greater than $B$ is less than $1/100$.
Then choose a fixed $0<r<1$ so small that
\eqref{ht:small-tail} makes the probability of any of
\begin{align*}
 &p\ne\widehat p,\quad
        |\sigma(Z_{\widehat p}-Z_p)|<r,\\
 &D_j\ne D_l,\quad |\sigma(D_j-D_l)|<r
       \qquad(1\le j<l\le3)
\end{align*}
less than $1/100$ in total. The equalities have deliberately been excluded
from these small-difference events. These estimates hold for every $H$
because they were integrated over a fresh independent embedding using
only the fixed marginal bounds.

Choose now
\begin{equation}\label{ht:mesh}
 0<\tau<\min\left(1,\frac{r^4}{128B^5}\right).
\end{equation}
At an embedding, encode each point by quantizing the real and imaginary
parts of both $Z_p,W_p$ in $[-B-1,B+1]$, in half-open intervals of length
at most $\tau$. Add an overflow code for points outside this box, using
an endpoint convention to make a partition. This is a finite alphabet
of size $M=M(B,\tau)$ independent of $t$ and of the field degree.
An acceptance test at $q$ asks whether the point code is attained by
some actual neighbor of $q$. The actual point on an edge always passes.

By \Cref{inc:entropy}, $D=o(\log t)$. Choose $K\to\infty$ with
$K=o(\log t)$ and $D/K\to0$. Apply \Cref{pred:acceptance} to independent
uniform embeddings as tests. It gives a fixed strict-prefix history map
with at most $M^{K-1}=t^{o(1)}$ values for which the predicted acceptance
failure at a fresh embedding is $o(1)$. The choice of history is made
after averaging the prediction failure over that fresh embedding; the
embedding itself is not fixed. The geometric error bounds just proved
remain valid for this selected history.

Use this $H$ in \eqref{ht:coupling}, with one common prediction
$\widehat p$ and three independently sampled neighbors of $p$.
Each pair $(\widehat p,q_j)$ has the prediction law in the acceptance
estimate, so all three tests pass outside probability $o(1)$ by a union
bound. Moreover
\begin{equation}\label{ht:collisions}
 \Prob(\widehat p=p)
   =\sum_{h:\mu(H=h)>0}
      \frac{\sum_{p:H(p)=h}\mu(p)^2}{\mu(H=h)}
   \le |H(P)|\max_p\mu(p)=t^{-1+o(1)}.
\end{equation}
For $j\ne l$, the conditional uniform neighbor law gives
$\Prob(q_j=q_l)\le1/\min_p d(p)=o(1)$.
Distinct centers have distinct $Z$-coordinates, so $q_j\ne q_l$ implies
$D_j\ne D_l$. The degree and atom estimates used here follow from
\Cref{inc:piece}. Combining all these estimates gives positive
probability that all three tests pass, all displayed quantities are
bounded by $B$, and
\begin{equation}\label{ht:separation}
 |\sigma(Z_{\widehat p}-Z_p)|\ge r,\qquad
 |\sigma(D_j-D_l)|\ge r\quad(j<l).
\end{equation}

Fix such an outcome and suppress $\sigma$ from the notation.
Successful acceptance provides, for each $j$, a genuine neighbor of
$q_j$ in the same code as $\widehat p$. Since the latter has coordinates
bounded by $B$, this is a nonoverflow code, and the two complex coordinate
errors are at most $\sqrt2\tau$. Comparing with that neighbor's exact
edge equation gives
\[
 |(Z_{\widehat p}-Z_{q_j})(W_{\widehat p}-W_{q_j})-1|
       \le8B\tau.
\]
Set $\Delta Z=Z_{\widehat p}-Z_p$ and
$\Delta W=W_{\widehat p}-W_p$. Subtracting the exact equation for
$(p,q_j)$ yields
\[
 \Delta Z/D_j+\Delta W D_j+\Delta Z\Delta W=E_j,
       \qquad |E_j|\le8B\tau.
\]
Subtract the equation with $j=1$ from those with $j=2,3$.
The coefficient matrix for $(\Delta Z,\Delta W)$ has determinant
\begin{equation}\label{ht:determinant}
 \det\begin{pmatrix}
   D_2^{-1}-D_1^{-1}&D_2-D_1\\
   D_3^{-1}-D_1^{-1}&D_3-D_1
 \end{pmatrix}
 =\frac{(D_1-D_2)(D_1-D_3)(D_2-D_3)}{D_1D_2D_3}.
\end{equation}
Its absolute value is at least $r^3/B^3$ by
\eqref{ht:separation}. Each matrix entry is at most $2B$, and each
right-hand side is at most $16B\tau$. Cramer's rule therefore gives
\[
 |\Delta Z|\le\frac{64B^5\tau}{r^3}<r/2,
\]
contradicting \eqref{ht:separation}. This excludes every bounded
subsequence of $S$ and proves the proposition.
\end{proof}

\subsection{Coordinate states on the original levels}

We now return to the original coordinates $z,w$, intervals $I_v$, and
thresholds $n_e$. By \Cref{ht:bounded}, $S\to\infty$.
At a level $(v,s)$ give each vertex $x$ a state
\[
 y_x(v,s)=(A_x(v,s),B_{1,x}(v,s)).
\]
At a nonarchimedean entry, $A_x$ is the equivalence class for
$|z_x-z_{x'}|_v\le e^{-s}$, and $B_{1,x}$ is that for
$|w_x-w_{x'}|_v\le e^s$. At an archimedean entry associated with $\sigma$,
$A_x$ is the half-open square-grid cell containing $\sigma(z_x)$ at
side length $e^{-s}$, and $B_{1,x}$ is the cell containing $\sigma(w_x)$
at side length $e^s$. For each of these two grids choose a fractional
shift in $[0,1)^2$ and use that same fractional shift, scaled with the
side length, at every level of the entry.

For a state $y=(a,b)$ its strict cross is
\begin{equation}\label{ht:strict-cross}
 R_y=\{(a',b'):\text{exactly one of }a'=a\text{ and }b'=b\text{ holds}\}.
\end{equation}
Call an edge \emph{consistent} at a level if $Y_e=1$ and exactly its
$A$-labels agree, or if $Y_e=0$ and exactly its $B_1$-labels agree.
Thus consistency implies that the endpoint states satisfy strict cross,
and also specifies which coordinate agrees. Define
\begin{equation}\label{ht:inconsistency}
 m(v,s)=\Prob_E\{e\text{ is inconsistent at }(v,s)\},\qquad
 m_e=\int\mathbf1_{\{e\text{ is inconsistent}\}}.
\end{equation}

\Needspace{5\baselineskip}
\begin{lemma}[State consistency]\label{ht:states}
The fractional shifts can be fixed so that
\begin{equation}\label{ht:inconsistency-budget}
 \int m=\E_E m_e=O(1),
\end{equation}
with an absolute constant. At each nonarchimedean entry every edge is
consistent outside finitely many threshold levels. If two vertices
have matching $A$-labels, then
\[
 |z_x-z_{x'}|_v\le c_v e^{-s};
\]
matching $B_1$-labels similarly implies
$|w_x-w_{x'}|_v\le c_v e^s$, where $c_v=1$ at nonarchimedean entries
and $c_v=\sqrt2$ at archimedean entries.
All the level functions can be chosen measurable.
\end{lemma}

\begin{proof}
At a nonarchimedean entry, $n_e>s$ means
$|d_e|_v<e^{-s}$ and $|d_e^{-1}|_v>e^s$, so exactly $A$ agrees.
For $n_e<s$ the two conclusions reverse. Only $s=n_e$ is excluded,
a null set for the level measure.

At infinity, put $r=|s-n_e|$. The smaller coordinate jump divided by
its grid scale is $e^{-r}$, so its mismatch probability under a uniform
shift is at most $\sqrt2e^{-r}$. The larger jump divided by its scale
is $e^r$; it cannot lie in the same cell once $r\ge\log\sqrt2$,
by the diameter of a half-open square. Consequently the expected
inconsistency indicator is at most
\[
 \mathbf1_{\{|s-n_e|\le\log\sqrt2\}}+
                       \sqrt2e^{-|s-n_e|}.
\]
Its integral over the full real line is bounded by an absolute constant,
and hence so is its integral over $I_v$. Sum over entries using total
archimedean weight one, then average over edges. There are only finitely
many archimedean shifts to choose at each index, so averaging fixes
shifts satisfying \eqref{ht:inconsistency-budget}.

The match implications follow directly from the equivalence classes and
the square diameter. For measurability, label nonarchimedean classes by
the first vertex in a fixed ordering of $P\cup Q$, and use integer cell
indices at infinity. All class comparisons are threshold comparisons
and all cell indices are measurable functions of $s$.
\end{proof}

We have obtained a space of tests of finite measure with minority
mass $S\to\infty$, while the integrated discrepancy between an edge's
threshold bit and its two coordinate agreements remains $O(1)$.
This makes that discrepancy negligible relative to $S$.
